\section{Related Work}

\paragraph{Weak test suites}
The problem of weak test suites has been well studied in the area of test-based automatic program repair (APR)~\cite{martinez2015automatic, qi2015analysis, smith2015cure, liu2021critical}.
For instance, 
Martinez et al.~\cite{martinez2015automatic} manually assessed 84 plausible patches generated on the Defects4J benchmark~\cite{just2014defects4j}, and found that only 11 of them are correct. 
Qi et al. ~\cite{qi2015analysis} found that most of the plausible patches 
are equivalent to a single modification that deletes the corresponding functionality, which is problematic. 
Some prior work proposed actionable strategies to mitigate this problem in APR~\cite{smith2015cure, xiong2018identifying, yu2019alleviating, yang2023large}. 
For example, 
Smith el al.~\cite{smith2015cure} found that using a test suite that is inaccessible to evaluated tools can help detect plausible but incorrect patches.
Xiong et al.~\cite{xiong2018identifying} proposed to incorporate execution behavior similarity under generated test inputs to detect incorrect patches that behave significantly differently from their oracle patches.
For test-based APR benchmarks, test suites are provided to the evaluated tools, so plausible but incorrect patches are also called overfitting patches.
Different from these studies, our work focuses on the issue-solving task, which aims to resolve issues based on issue statements rather than test suites.
For issue-solving benchmarks, the evaluated tools cannot access test suites during evaluation.
Thus, prior conclusions may not apply to SWE-bench.
Recently, Aleithan et al.~\cite{aleithan2024swe} also noticed the weak test suite problem in SWE-bench.
However, first, their study is conducted on the patches generated by only one issue-solving tool on SWE-bench full, which is shown to contain many low-quality instances.
In contrast, our study covers three state-of-the-art tools and focuses on the human-filtered subset SWE-bench Verified.
Second, they focus on manually classifying plausible patches, while we dig deeper into the patterns and types of plausible but incorrect patches with our novel technique \appname.



\paragraph{Automated test generation}
Various automated test generation approaches are proposed to alleviate testing efforts and help find bugs\cite{fraser2011evosuite, lukasczyk2022pynguin, pacheco2007randoop, agitar2023}.
\emph{Traditional test generation} approaches utilize some predefined rules to generate tests and can be mainly categorized into search-based~\cite{fraser2011evosuite, lukasczyk2022pynguin}, randomization-based~\cite{pacheco2007randoop}, and constraints-based~\cite{tillmann2008pex, chang2016constraint} approaches.
\emph{Deep learning-based test generation} approaches train deep learning models to generate natural and human-understandable test cases~\cite{tufano2020unit, alagarsamy2024a3test, dinella2022toga, rao2023cat}.
Recently, researchers have proposed \emph{LLM-based test generation} approaches~\cite{chen2024chatunitest, ryan2024code, pizzorno2024coverup, wang2024hits}.
For example, CoverUp iteratively instructs LLMs to generate Python regression tests, improve coverage, and fix errors with detailed coverage information.
Liu et al. ~\cite{liu2024your} proposed to augment the test suites in the code generation benchmarks HumanEval~\cite{chen2021evaluating} and MBPP~\cite{austin2021program} based on LLMs and type-aware mutation.
These approaches aim to generate regression tests to cover as much of the code as possible, while \appname targets differentiating two patches, which requires more targeted test generation.
Differential testing aims at generating tests to expose different behaviors between two versions of a program ~\cite{evans2007differential, li2023nuances, liu2024llm, etemadi2024mokav}. 
Mokav~\cite{etemadi2024mokav} uses execution feedback to iteratively guide an LLM in generating tests to differentiate different solutions of programming competition tasks.
Existing LLM-based differential testing techniques focus on function-level programs, while \appname can generate differentiating tests for large and complex projects.
In addition, \appname targets differential patch testing, and leverages specially designed methods to identify appropriate target functions and to construct useful contextual information for revealing meaningful behavioral differences, which are important for differential patch testing but are not considered by prior work.












\paragraph{Software engineering agents}
LLM-based agents promise to help automate various software engineering tasks~\cite{AgenticAISE2025}.
Agents for various tasks have been presented, such as fault localization~\cite{kang2024evaluating,xu2025flexfl,jiang2025cosil}, issue solving~\cite{Yang2024a,zhang2024autocoderover,xia2024agentless}, automated program repair~\cite{icse2025-RepairAgent,Cheng2025}, repository setup~\cite{issta2025_ExecutionAgent,Eliseeva2025,Milliken2025}, and generating issue-reproducing tests~\cite{Muendler2024,Issue2Test_arXiv2025}.
Researchers have started to study the behavior such agents to better understand their strengths and weaknesses~\cite{ase2025_agent_study}.
Our work complements such efforts and will help improve future software engineering agents by critically analyzing the patches produced by the agents.
